% !TeX program = pdflatex
% =============================================================================
% Lektionsplanung (Doppellektion, 90 Minuten): s-t-Diagramm
% Kantonsschule KFR -- Physik
% Internes Planungsdokument fuer die Lehrperson, nicht schuelerseitig.
% =============================================================================
\documentclass[11pt,a4paper]{article}

\usepackage[T1]{fontenc}
\usepackage[utf8]{inputenc}
\usepackage[ngerman]{babel}
\usepackage{lmodern}
\usepackage[a4paper,margin=1.6cm,headheight=14pt]{geometry}
\usepackage{array,booktabs,tabularx,multirow,longtable}
\usepackage{enumitem}
\usepackage{siunitx}
\sisetup{per-mode=fraction,fraction-command=\frac}
\usepackage{xcolor}
\usepackage{colortbl}
\usepackage[most]{tcolorbox}
\usepackage{lastpage}
\usepackage{fancyhdr}
\usepackage{setspace}
\setlength{\parindent}{0pt}
\setlength{\parskip}{0.5em}
\setstretch{1.05}

% -----------------------------
% Farben (analog Corporate-Farbschema Physik KFR)
% -----------------------------
\definecolor{titleblue}{HTML}{183A6B}
\definecolor{accentblue}{HTML}{2E6F9E}
\definecolor{lightblue}{HTML}{EAF4FB}
\definecolor{lightgray}{HTML}{F4F6F8}
\definecolor{darkgray}{HTML}{4A4A4A}
\definecolor{groupteal}{HTML}{1B7F72}
\definecolor{lightgroupteal}{HTML}{E8F5F3}
\definecolor{theorygreen}{HTML}{2E7D4F}
\definecolor{lighttheorygreen}{HTML}{EAF7EF}

% -----------------------------
% Titelbanner
% -----------------------------
\newtcolorbox{titlebox}{
  enhanced, colback=titleblue, colframe=titleblue, boxrule=0pt, arc=3mm,
  left=3mm,right=3mm,top=3mm,bottom=3mm
}
\newtcolorbox{infobox}{
  enhanced, breakable, colback=lightblue, colframe=accentblue, boxrule=0.7pt,
  arc=2mm, left=6mm,right=6mm,top=3.5mm,bottom=3.5mm
}
\newtcolorbox{hinweisbox}[1][Hinweis]{
  enhanced, breakable, colback=lighttheorygreen, colframe=theorygreen,
  title={#1}, fonttitle=\bfseries, boxrule=0.7pt, arc=2mm,
  left=6mm,right=6mm,top=3.5mm,bottom=3.5mm
}

\pagestyle{fancy}
\fancyhf{}
\lhead{\textcolor{darkgray}{Lektionsplanung: s-t-Diagramm}}
\rhead{\textcolor{darkgray}{Physik, 4a}}
\cfoot{\textcolor{darkgray}{Seite \thepage\ von \pageref{LastPage}}}
\renewcommand{\headrulewidth}{0.4pt}

% -----------------------------
% Sozialform-Kuerzel farbig hervorheben
% -----------------------------
\newcommand{\EA}{\textcolor{accentblue}{\textbf{EA}}}
\newcommand{\PA}{\textcolor{groupteal}{\textbf{PA}}}
\newcommand{\PL}{\textcolor{titleblue}{\textbf{Plenum}}}
\newcommand{\LV}{\textcolor{darkgray}{\textbf{LV}}}

\setlength{\tabcolsep}{4pt}
\newcolumntype{Z}{>{\raggedright\arraybackslash}p{1.1cm}}
\newcolumntype{I}{>{\raggedright\arraybackslash}p{4.4cm}}
\newcolumntype{L}{>{\raggedright\arraybackslash}p{4.4cm}}
\newcolumntype{U}{>{\raggedright\arraybackslash}p{4.2cm}}
\newcolumntype{F}{>{\centering\arraybackslash}p{1.5cm}}

\begin{document}

\begin{titlebox}
{\color{white}\sffamily\bfseries\Large Lektionsplanung (Doppellektion): s-t-Diagramm}
\end{titlebox}

\noindent\textbf{Klasse:}~4a \hfill \textbf{Datum:}~21.08.2026 \hfill
\textbf{Thema:}~s-t-Diagramm \hfill \textbf{Dauer:}~90 Minuten
\medskip

\noindent\textbf{Lehrperson:}~Loris Keller \hfill
\textbf{Material:}~Arbeitsblatt ,,s-t-Diagramm`` (Klassensatz), Wandtafel/Whiteboard,
Beamer oder Dokumentenkamera (fuer Bildfahrplan, SBB-App-Screenshot,
Kilometerstein- und Blitzerfoto)

\begin{infobox}
\textbf{Lernzielfokus dieser Doppellektion}\\
Diese Doppellektion deckt \textbf{LZ1 bis LZ7} des Arbeitsblatts vollstaendig
ab (Weg vs.\ Ortsänderung, vertikale Linie unmöglich, gleichförmige
Bewegung, Sekantensteigung als Durchschnittsgeschwindigkeit, Vorzeichen der
Steigung, Bewegungsbeschreibung $\leftrightarrow$ Diagramm). \textbf{LZ8 und
LZ9} (Tangente/Momentangeschwindigkeit, Grenzwert $\Delta t \to 0$) sind
bewusst \textbf{nicht} Teil dieser Doppellektion. Das Arbeitsblatt ist mit
ca.\ 90 Punkten und 25 Seiten deutlich zu umfangreich fuer 90 Minuten. Sie
werden in der folgenden Lektion behandelt (siehe Abschnitt ,,Ausblick`` am
Ende dieses Dokuments). Die Zusatzaufgaben (Aufholjagd, Zwei Züge treffen
sich) sowie die Übung ,,Bildfahrplan mit drei Zügen`` sind als Vertiefung
bzw.\ Hausaufgabe vorgesehen, nicht als Unterrichtsstoff dieser Lektion.
\end{infobox}

\begin{hinweisbox}[Sozialform-Kuerzel]
\EA{} Einzelarbeit \quad \PA{} Partnerarbeit \quad \PL{} Plenum
(Klassengespräch) \quad \LV{} Lehrervortrag (Input der Lehrperson)
\end{hinweisbox}

\section*{1.\ Lektion (45 Minuten)}

\begin{longtable}{Z|I|L|U|F}
\toprule
\textbf{Zeit} & \textbf{Phase / Inhalt} & \textbf{Lehrperson} & \textbf{SuS} & \textbf{Sozial\-form}\\
\midrule
\endhead
0--8' &
Einstieg: Hook ,,Woher weiss die SBB, wann der Zug kommt?`` &
Zeigt Bildfahrplanfoto, laesst Klasse zunaechst spekulieren (keine
Aufloesung), zeigt danach SBB-App-Screenshot und stellt Leitfrage der
Lektion. &
Aeussern Vermutungen zum Bildfahrplan, hoeren zu, formulieren eigene
Fragen. &
\PL{} \\
8--14' &
Theorie: Ort $s$, Zeit $t$, Aufbau des $s$-$t$-Diagramms &
Fuehrt Begriffe Ort/Zeit/Nullpunkt ein, zeichnet Achsen an der Tafel,
liest gemeinsam einen Punkt $(t,s(t))$ ab. &
Uebernehmen Definitionen ins Heft/Arbeitsblatt, lesen den Beispielpunkt
mit. &
\LV{} \\
14--17' &
Beispiel: Auto auf der Autobahn (Kilometerstein) &
Zeigt Kilometersteinfoto, verknuepft $s=\SI{2.6}{\kilo\meter}$ mit der Anzeige auf
dem Foto, entwickelt das zugehoerige Diagramm an der Tafel. &
Ordnen die Fotobeschriftung der Grösse $s$ zu, verfolgen die
Diagrammkonstruktion. &
\PL{} \\
17--22' &
Theorie: Zurückgelegter Weg vs.\ Ortsänderung &
Stellt die beiden Bewegungen mit gleicher Ortsänderung, aber
unterschiedlichem Weg gegenueber; fragt nach, worin der Unterschied liegt. &
Vergleichen die beiden Diagramme, formulieren den Unterschied in eigenen
Worten. &
\LV{} / \PL{} \\
22--32' &
Übung: Weg vs.\ Ortsänderung (Velofahrerin) &
Gibt Arbeitsauftrag, steht fuer Rueckfragen zur Verfuegung; bespricht
anschliessend die Loesung an der Tafel. &
Bearbeiten die Aufgabe selbststaendig (7'), vergleichen danach mit der
Musterloesung (3'). &
\EA{} \,\textrightarrow\, \PL{} \\
32--36' &
Theorie: Die Zeit läuft nur vorwärts &
Zeigt beide Beispielgraphen (erlaubt/nicht erlaubt), fragt die Klasse,
welcher physikalisch unmoeglich ist und warum. &
Beurteilen die beiden Graphen, begruenden ihre Einschaetzung im
Klassengespräch. &
\PL{} \\
36--45' &
Theorie: Lineare Bewegung: Geschwindigkeit als Steigung &
Leitet $v=\Delta s/\Delta t$ her, betont den Unterschied zu einer
einzelnen abgelesenen Koordinate, zeigt zwei verschieden steile Geraden. &
Verfolgen die Herleitung, uebernehmen die Formel samt Einheit. &
\LV{} \\
\bottomrule
\end{longtable}

\textit{Kurze Pause (5 Minuten) zwischen den beiden Lektionen.}

\section*{2.\ Lektion (45 Minuten)}

\begin{longtable}{Z|I|L|U|F}
\toprule
\textbf{Zeit} & \textbf{Phase / Inhalt} & \textbf{Lehrperson} & \textbf{SuS} & \textbf{Sozial\-form}\\
\midrule
\endhead
0--4' &
Theorie: Einheiten umrechnen m/s $\leftrightarrow$ km/h &
Leitet den Umrechnungsfaktor $3.6$ an der Tafel her (Autotacho als
Anwendungsbeispiel). &
Verfolgen die Herleitung, notieren die Umrechnungsformel. &
\LV{} \\
4--12' &
Übung: Steigung ablesen und Einheiten umrechnen &
Gibt Arbeitsauftrag; bespricht danach kurz das Ergebnis (insbesondere die
Umrechnung ohne vorgegebene Formel, Teilaufgabe b). &
Lesen Koordinaten aus dem Diagramm ab, berechnen $v$ und rechnen
selbststaendig in km/h um (6'), vergleichen im Plenum (2'). &
\EA{} \,\textrightarrow\, \PL{} \\
12--22' &
Übung: Blitzer mit Induktionsschleifen &
Gibt Arbeitsauftrag fuer Partnerarbeit, geht herum und unterstuetzt bei
Rueckfragen; moderiert die Schlussbesprechung inkl.\ Diagrammvergleich
(eigene Geschwindigkeit vs.\ die Gerade des Tempolimits). &
Berechnen zu zweit Geschwindigkeit und Blitzerentscheidung, zeichnen
gemeinsam beide Geraden ins Diagramm. &
\PA{} \,\textrightarrow\, \PL{} \\
22--27' &
Theorie: Betrag einer Zahl + die drei Grundbewegungen &
Erklaert den Betrag anhand der Zahlengeraden, ordnet Steigungsvorzeichen
den drei Grundbewegungen Anhalten, vorwaerts und rueckwaerts zu. &
Uebertragen die Zuordnung Vorzeichen $\leftrightarrow$ Richtung ins
Arbeitsblatt. &
\LV{} \\
27--31' &
Kurzaufgabe: Vertikale Linie &
Laesst die Aufgabe zunaechst einzeln loesen, ruft danach eine Erklaerung
ab und ergaenzt bei Bedarf. &
Begruenden schriftlich, weshalb der senkrechte Abschnitt physikalisch
keinen Sinn ergibt. &
\EA{} \,\textrightarrow\, \PL{} \\
31--36' &
Theorie: Bewegungsabschnitte im kombinierten Diagramm + Beispiel
Raststätte &
Zeigt das mehrteilige Autobahnbeispiel, hebt hervor, dass gleiche
Geschwindigkeit an unterschiedlichen Stellen gleiche Steigung bedeutet. &
Ordnen die einzelnen Abschnitte des Raststättenbeispiels den drei
Grundbewegungen zu. &
\PL{} \\
36--43' &
Partnerarbeit: Bewegungsgeschichte zeichnen &
Erklaert den Ablauf (zuerst einzeln denken, dann beschreiben:
Partner/-in zeichnet nur nach Wortbeschreibung), sammelt am Ende ein bis
zwei Beispiele mit Missverstaendnissen fuer die Klasse ein. &
Erfinden allein eine Bewegungsgeschichte und zeichnen Diagramm~1;
beschreiben sie danach nur muendlich, waehrend die Partnerin/der Partner
Diagramm~2 zeichnet; vergleichen beide Diagramme. &
\PA{} \\
43--45' &
Abschluss, Sicherung und Hausaufgabe &
Fasst die Kernaussage der Doppellektion zusammen ($v=\Delta s/\Delta t$
als Steigung, Vorzeichen = Richtung), gibt Hausaufgabe und Ausblick auf
naechste Lektion (Momentangeschwindigkeit). &
Hoeren zu, notieren die Hausaufgabe. &
\PL{} \\
\bottomrule
\end{longtable}

\begin{hinweisbox}[Hausaufgabe]
Übung ,,Bildfahrplan mit drei Zügen`` (Wiederholung Sekantensteigung,
Begegnen/Überholen) sowie (falls Zeit reicht) die beiden
Zusatzaufgaben (,,Aufholjagd beim Radrennen``, ,,Zwei Züge treffen sich``)
als Vorbereitung und Übungsmaterial. Keine dieser drei Übungen setzt bereits
Momentangeschwindigkeit (LZ8/LZ9) voraus.
\end{hinweisbox}

\begin{hinweisbox}[Ausblick auf die Folgelektion]
Die naechste Lektion (Einzellektion, 45 Minuten) deckt \textbf{LZ8} und
\textbf{LZ9} ab: Theoriebox ,,Durchschnitts- und Momentangeschwindigkeit``
(Sekante vs.\ Tangente, Grenzwert $\Delta t \to 0$), anschliessend die
Übungen ,,Momentan- vs.\ Durchschnittsgeschwindigkeit`` (Aufzug,
stückweise linear) und ,,Momentan- und Durchschnittsgeschwindigkeit an
einer Kurve`` (Kugel auf Rampe, echtes Tangentenzeichnen an einer
gekrümmten Kurve). Beide Lernziele werden im aktuellen Arbeitsblatt nur
theoretisch behauptet, aber noch nicht durch eine Aufgabe aktiv
eingefordert (LZ3, gleichförmige Bewegung erklären, ebenso); das sollte
bei der Ueberarbeitung des Arbeitsblatts ergaenzt werden.
\end{hinweisbox}

\end{document}
